\par\medskip\noindent
2.3. \underline{Relación con la ``dualidad de Cartier''.}
La descripción dada en 2.2 de los acoplamientos definidos por una
biextensión, mediante el teorema de isomorfismo VII 3.6.5 y el
homomorfismo (2.2.1), no es manifiestamente la definición utilizada
antiguamente para los esquemas abelianos, cuando todavía no existía la
noción misma de biextensión. Aquella definición procede por medio de
la dualidad de Cartier; vamos a relacionarla con la actual.

Consideremos de nuevo el homomorfismo
\begin{equation}\tag{2.3.1}
 \Biext^1(P,Q;G)
 \longrightarrow
 \Hom\bigl(P,\uExt^1(Q,G)\bigr)
\end{equation}
procedente de la sucesión exacta (1.4). Es, de hecho, un isomorfismo
cuando
\begin{equation}\tag{2.3.2}
 \uHom(Q,G)=0.
\end{equation}

%%%%%%%%%%  scan idx 253  |  folio 242  %%%%%%%%%%

Esta condición se cumple en particular cuando $Q$ es un esquema
abeliano sobre un esquema $S$ y $G$ es el grupo multiplicativo
$\underline G_m$ (VII 1.3.8).

Existe también un homomorfismo evidente de ``restricción a $nQ$''
\[
 \uExt^1(Q,G)\longrightarrow\uExt^1(nQ,G),
\]
y por tanto un homomorfismo compuesto
\begin{equation}\tag{2.3.4}
\begin{split}
 \Hom\bigl(P,\uExt^1(Q,G)\bigr)
 &\longrightarrow
 \Hom\bigl(P,\uExt^1(nQ,G)\bigr)\\
 &\longrightarrow
 \Hom\bigl({}_nP,{}_n\uExt^1(nQ,G)\bigr),
\end{split}
\end{equation}
donde la segunda flecha es la restricción a $\,{}_nP$. Componiendo con
(2.3.1), obtenemos un homomorfismo
\begin{equation}\tag{2.3.5}
 \Biext^1(P,Q;G)
 \longrightarrow
 \Hom\bigl({}_nP,{}_n\uExt^1(nQ,G)\bigr).
\end{equation}

Por otra parte, la sucesión exacta
\begin{equation}\tag{2.3.6}
 0\longrightarrow{}_nQ\longrightarrow Q
 \xrightarrow{\ n\ }nQ\longrightarrow0
\end{equation}
da lugar a una sucesión exacta larga de los $\uExt^i(-,G)$, y por
tanto a la sucesión exacta corta
\begin{equation}\tag{2.3.7}
 0\longrightarrow\uHom(Q,G)_n
 \longrightarrow\uHom({}_nQ,G)
 \xrightarrow{\ \partial\ }
 {}_n\uExt^1(nQ,G)
 \longrightarrow0,
\end{equation}
donde, como de costumbre, $E_n$ denota $E/nE$. Retenemos de esta
sucesión el homomorfismo canónico suprayectivo
\begin{equation}\tag{2.3.8}
 \uHom({}_nQ,G)
 \xrightarrow{\hspace{3em}\partial\hspace{3em}}
 {}_n\uExt^1(nQ,G),
\end{equation}
que es un isomorfismo si $\uHom(Q,G)_n=0$, y con mayor razón si se
cumple (2.3.2). De (2.3.8) deducimos un homomorfismo
\begin{equation}\tag{2.3.9}
\begin{split}
 \Hom\bigl({}_nP,\uHom({}_nQ,G)\bigr)
 &\simeq
 \Hom({}_nP\otimes{}_nQ,G)\\
 &\longrightarrow
 \Hom\bigl({}_nP,{}_n\uExt^1(nQ,G)\bigr),
\end{split}
\end{equation}
que es un isomorfismo cuando se cumple (2.3.2).

Todos los homomorfismos considerados son funtoriales en $P,Q,G$. Los
homomorfismos (2.3.5), (2.3.9) y (2.2.2) forman el diagrama

%%%%%%%%%%  scan idx 254  |  folio 243  %%%%%%%%%%

\begin{equation}\tag{2.3.10}
\begin{tikzcd}[column sep=large,row sep=large]
 &\Biext^1(P,Q;G)
   \arrow[ddl,"\varphi^n"']
   \arrow[ddr,"(2.3.5)"]&\\
 &-&\\
 \Hom({}_nP\otimes{}_nQ,G)
   \arrow[rr,"(2.3.9)"]&&
 \Hom\bigl({}_nP,{}_n\uExt^1(\sisi{nQnG}{nQ,G})\bigr),
\end{tikzcd}
\end{equation}
donde la flecha horizontal (2.3.9) es un isomorfismo si se cumple
(2.3.2).

\medskip
\noindent\underline{Proposición 2.3.11.}
\underline{El diagrama (2.3.10) es anticonmutativo.}

Cuando $\uHom(Q,G)=0$, este enunciado recupera, por tanto, una
caracterización del homomorfismo $\varphi^n$ de (2.2.2) en términos
de los homomorfismos (2.3.5) y (2.3.9). Es la caracterización
utilizada antiguamente en la teoría de los esquemas abelianos, donde
además $nQ=Q$.
